\documentclass[aspectratio=169]{beamer}

\usetheme{default}

\title{Theme Gallery}
\subtitle{All 12 built-in themes}
\author{Slitex}
\date{\today}

\begin{document}

\begin{frame}
  \titlepage
\end{frame}

\begin{frame}{Available Themes}
  Change the theme with \texttt{\textbackslash usetheme\{Name\}}.

  \bigskip

  \begin{columns}
    \begin{column}{0.48\textwidth}
      \begin{itemize}
        \item \textbf{default}
        \item \textbf{metropolis}
        \item \textbf{madrid}
        \item \textbf{warsaw}
        \item \textbf{frankfurt}
        \item \textbf{copenhagen}
      \end{itemize}
    \end{column}
    \begin{column}{0.48\textwidth}
      \begin{itemize}
        \item \textbf{berlin}
        \item \textbf{boadilla}
        \item \textbf{cambridgeus}
        \item \textbf{darmstadt}
        \item \textbf{annarbor}
        \item \textbf{berkeley}
      \end{itemize}
    \end{column}
  \end{columns}
\end{frame}

\begin{frame}[plain]{Plain Frame}
  A frame with the \texttt{[plain]} option has no header or footer decorations.

  Useful for title cards or full-screen images.
\end{frame}

\begin{frame}{Frame with Subtitle}
  \framesubtitle{This is the frame subtitle}

  The \texttt{\textbackslash framesubtitle} command adds a subtitle line below the frame title.

  You can also pass it as the second brace argument: \texttt{\textbackslash begin\{frame\}\{Title\}\{Subtitle\}}.
\end{frame}

\end{document}
